% Transcription — fonds Alexandre Grothendieck, shelfmark 15, batch 4 (pages 61-80).
% Established from the facsimile archives/batches/15/batch-04.pdf (a mirror of
% https://grothendieck.umontpellier.fr/15.pdf), per the skill
% transcribe-grothendieck. The batch file carries no cover sheet: PDF page N
% is page 60+N of the folder (the Montpellier footer prints « 61 » ... « 80 »),
% and the archivists' pencil numbers bottom-left agree on every transcribed
% leaf (62, 65, 67, 69, 71, 73, 75, 76, 77, 80 all legible).
%
% Pass: Opus 5.5 (claude-opus-5-5), 2026-09-23 — first pass, unchecked against the pages by a human.
%
% Ten of the twenty pages are transcribed. Absent: pages 61, 63, 66, 68, 70,
% 72, 74, 78 — leaves of an unrelated typescript (a course on analytic
% functions of several variables: « Domaines d'existence maximaux »,
% « Dépendance des conditions initiales », « Résolution des équations
% analytiques », typed pp. 15-35, several scanned upside down), nothing in his
% hand; they are the versos of the handwritten sheets. Pages 64 and 79,
% blank. Page 76 is a leaf of the same typescript, upside down, but carries
% jottings in his hand unrelated to the typescript (Legendre symbol, p ≡ 1
% mod 4, Q(i)), which bear on page 69; it is transcribed for them alone.
%
% What the batch is. Two pieces.
%   (1) Pages 62-77, handwritten, fast, not paginated by him: the theory of
%   « catégories de Deuring » (A-linear categories locally equivalent to
%   P(o/A), o = End(X) a maximal order), their Brandt groupoid and Picard
%   category.
%   62     continues page 60 (batch 3) mid-sentence: « (tout objet » / de C
%          est de la forme P ⊗_o X (o = End X), P projective of finite type,
%          when the tensor products are « raisonnables »; C is then a full
%          subcategory of P_d(o) containing the unit object. A « jolie »
%          category is the inductive limit of its « strictement monogènes »
%          full subcategories, each equivalent to a subcategory of some P_d(o).
%   65     « C catégorie de Deuring (Picard) », C ≃ P(o/A). Brdt C acts on C,
%          Brdt(C) = π_0 acts on Is(C); quotient = classes of algebras
%          apparentées; stabiliser = Autext(o) if Brdt(A)=0 (sic). Exact
%          sequence 0 → o^{Z*} → o^* → Aut_A o → Brdt(C) → Is C → Is O → 0,
%          then in boxed form π_1(B,1) → π_1(C,X) → π_1(O,o) → π_0(B,1) →
%          π_0(C,X) → π_0(O,o) → 1 for the Brandt category B, the Deuring
%          category C and the category O of rings apparentés. « On a envie de
%          dire que C est fibrée sur O, de fibre B. » Pic(A) → B, π_1(B,1) =
%          A^*, Pic(A) → π_0(B,1) → Ram(C/A) → 0, Ram ≃ ∐_p Z/e_p Z.
%   67     same material restated: Eqv(C) gr-category, Brd(C) = π_0 Eqv
%          acting on classes of C; Pic(A) → Eqv(C) fully faithful, L ↦ L ⊗_A X;
%          · → Pic(A) → Brdt(C) → Brdt(C)/Pic(A) → 0; in the Brandt case
%          Brdt C/Pic(A) ≃ Div(C)/Div(A) = ∐_{p ramified} Z/e_p Z.
%   69     A local Dedekind: Is(C)=1, Brd(C)=Z/eZ, Autext(o)=Z/eZ (local
%          obstruction to automorphisms being inner); when they are locally
%          inner, the obstruction lies in Pic(A). « Cas des corps gauches de
%          Deuring » (ramified at p and ∞): 0 → Z/2Z → o^* → Aut(o) → Z/2Z →
%          π_0(C,X) → π_0(O,o) → 1, with o^* = Z/2Z except j = 0, 12^3,
%          Frobenius, j ∈ F_p, E ↦ E^(p). Ends « Question : o^* est-il tjs
%          égal » — the sentence stops there.
%   71     O the category of algebras apparentées, K total ring of fractions
%          semi-local, E ≃ o_K; condition « tout autom. de E est intérieur »;
%          three 2-categories B, B_1 (bimodules loc. free of rank 1,
%          composition ⊗_{o'}), B_2 (Brandt 2-groupoid of suborders of E and
%          bimodules inside E), 2-functors B_2 → B_1 → B.
%   73-75  condition (C): L M = M L for special (o,o)-bimodules in E; hence a
%          strict commutativity constraint on Eqv(C) ≃ Hom_{B_1}(o,o), made
%          explicit via L ≃ L' ⊂ E. Proposition (p. 75): it does not depend on
%          X ∈ Ob C; proof by L ↦ P L P^{-1}, « OK. »
%   77     « Conclusion »: the Eqv(C) are strict Picard categories, all
%          canonically equivalent, hence one Picard category B (« catégorie
%          de Picard de Brandt »); Eqv(C,C') is a B-torsor-category; B_0 →
%          torsor-categories under B is an equivalence of 2-categories; base
%          change A → A'.
%   (2) Page 80: first leaf of a typescript, « Localisation pour les
%   variétés abéliennes. », § 1, with handwritten additions: C_{A'} with
%   Hom ⊗_A A', the fibred category C̲ over Spec(A') ⊂ S, « localisable sur
%   A », consequences a), b). It breaks off in the middle of b) (« le
%   foncteur P ↦ P ⊗_R X »), continued on page 81 (batch 5).
%
% Notation fixed for the batch. C (plain) for the Deuring category — his C
% has a loop on pages 65-77 and is plain on 62 and 80; o (lower case) for
% End(X), a maximal order; \mathcal{O} for the category of rings apparentés
% (his curly O); \mathcal{B}, \mathcal{B}_0, \mathcal{B}_1, \mathcal{B}_2 for
% the Brandt category and the 2-categories (he double-underlines B_0, B_1,
% B_2); \mathbb{P}(o/A), \mathbb{P}_d(o) for his double-struck P (categories
% of projective modules). Brdt, Brd (both spellings are his), Eqv (which he
% underlines), Is, Pic (double-underlined), Ram, Autext are kept as written.
% « t.f. », « loc. », « ens. », « autom. », « ½ local » (semi-local, per
% references/hand.md §2) are kept.
%
% Legibility. Pages 62, 71 and the margins of 65 and 73 are fast and heavily
% reworked; the formulas and the named objects carry, the connective prose
% often does not. Nothing on these leaves is dated.
%
% The dating is the inventory's own for the folder, copied verbatim. Its
% group, [10-18] « Théorie des motifs », is dated [à partir de 1958-à partir
% de 1982].
\documentclass[11pt,a4paper]{article}
% Le préambule commun à toutes les transcriptions du fonds Grothendieck.
%
% Il tient en une idée : **l'incertitude se compose, elle ne se résout pas.**
% Une transcription qui lisse un mot douteux a détruit la seule chose pour
% laquelle on la faisait — un lecteur doit pouvoir savoir, à chaque endroit,
% ce qui était sur le feuillet et ce qui a été ajouté. Les six macros qui
% suivent sont donc l'appareil critique, et non de la décoration.
%
% Les mêmes commandes sont reconnues par `scripts/render.mjs`, qui produit la
% vue de lecture du site. Une macro ajoutée ici doit l'être aussi là-bas,
% faute de quoi le rendu échouera bruyamment — ce qui est voulu : un rendu
% silencieusement faux est pire qu'un rendu absent.

\ProvidesPackage{grothendieck}[2026/08/08 Transcriptions du fonds Grothendieck]

\usepackage{amsmath,amssymb,amsthm}
\usepackage{mathrsfs}
\usepackage{stmaryrd}
% Les diagrammes commutatifs sont partout dans ce fonds : c'est la langue
% même de la géométrie algébrique de Grothendieck, et une transcription qui
% les rendrait en prose ne transcrirait rien.
\usepackage{tikz-cd}
% Français par défaut — c'est la langue des feuillets — mais l'anglais est
% chargé aussi : la traduction et le résumé sont des documents anglais, et la
% typographie française (espaces avant les deux-points) y serait une faute.
% Ces documents basculent par \selectlanguage{english} en tête de corps.
\usepackage[english,french]{babel}
\usepackage{fontspec}
\usepackage{geometry}
\geometry{a4paper,margin=25mm,marginparwidth=28mm}
\usepackage{marginnote}
\usepackage{xcolor}
% ulem, not soul. `soul`'s \st and \ul are built on a character-by-character
% scan that XeLaTeX's font machinery breaks: the document fails with an
% undefined control sequence rather than degrading. ulem does the same two jobs
% and is Unicode-engine safe. `normalem` stops it hijacking \emph, which the
% transcriptions use for Grothendieck's own emphasis.
\usepackage[normalem]{ulem}
\usepackage{hyperref}
\hypersetup{colorlinks=true,linkcolor=black,urlcolor=black,pdfborder={0 0 0}}

% --- Métadonnées du lot -----------------------------------------------------
% Reprises telles quelles par le script de rendu, qui les affiche en tête de
% la vue de lecture. Elles fixent ce que le fichier prétend couvrir : sans
% elles, une transcription flotte sans point d'attache dans un fonds de seize
% mille feuillets.

\newcommand{\folder}[1]{\gdef\grothFolder{#1}}
\newcommand{\batch}[1]{\gdef\grothBatch{#1}}
\newcommand{\leaves}[2]{\gdef\grothFirst{#1}\gdef\grothLast{#2}}
\newcommand{\foldertitle}[1]{\gdef\grothFoldertitle{#1}}
\gdef\grothFolder{?}\gdef\grothBatch{?}\gdef\grothFirst{?}\gdef\grothLast{?}\gdef\grothFoldertitle{}

% --- Le résumé --------------------------------------------------------------
%
% Il ouvre la lecture modernisée, sous le titre « Résumé », et il est composé
% en petit. Ce n'est pas un ornement : c'est le seul passage du document qui
% ne soit pas des mathématiques, et un lecteur doit pouvoir le reconnaître
% avant de l'avoir lu — puis le sauter s'il connaît déjà le sujet, ou s'y
% arrêter s'il ne le connaît pas. Le filet qui le suit marque la reprise du
% corps.

\newenvironment{resume}
  {\small\setlength{\parskip}{0.45em}}
  {\par\medskip\hrule\medskip}

% --- Mots-clés ---------------------------------------------------------------
%
% En anglais, en clôture du résumé de la lecture modernisée : le vocabulaire
% moderne sous lequel chercher ce que les feuillets construisent. Le manifeste
% du site (`npm run manifest`) les extrait pour étiqueter la cote entière —
% cette ligne est la source unique des tags, il n'y a pas de fichier à côté.
\newcommand{\keywords}[1]{\par\smallskip\noindent
  {\footnotesize\textsc{Keywords} --- \textit{#1}}\par}

% --- L'appareil critique ----------------------------------------------------

% Le numéro de feuillet, en marge. C'est l'ancre qui permet de régler un
% désaccord sur une formule en regardant *un* feuillet plutôt qu'une chemise
% de six cents. Le site s'en sert aussi pour tourner les pages du fac-similé
% quand on fait défiler la transcription.
\newcommand{\leaf}[1]{%
  \marginnote{\footnotesize\color{gray!70}\textsf{#1}}%
  \label{leaf:#1}\ignorespaces}

% Illisible : on ne devine pas. Un mot inventé qui se lit comme les autres est
% le pire résultat possible.
\newcommand{\ill}{\textcolor{red!60!black}{[\,\ldots\,]}}

% Lecture douteuse : le mot est donné, et signalé comme incertain.
\newcommand{\uncertain}[1]{\uline{#1}}

% Ajout de l'éditeur — une lettre manquante, un mot sous-entendu.
\newcommand{\add}[1]{\textcolor{blue!50!black}{[#1]}}

% Biffé par Grothendieck lui-même. Ce qu'il a rayé fait partie du document :
% une rature dit souvent où la pensée a changé de direction.
\newcommand{\struck}[1]{\sout{#1}}

% Note du transcripteur, sur l'état matériel du feuillet.
\newcommand{\note}[1]{\par\smallskip{\small\color{gray!60!black}\textit{[#1]}}\par\smallskip}

% Note marginale de Grothendieck, distincte du corps du texte.
\newcommand{\marginal}[1]{\par\smallskip\noindent\hspace{1em}%
  {\small\color{gray!60!black}#1}\par\smallskip}

% --- Titre ------------------------------------------------------------------

\AtBeginDocument{%
  \begin{center}
    {\large\bfseries Fonds Alexandre Grothendieck --- cote n\textsuperscript{o}~\grothFolder}\\[2pt]
    {\itshape\grothFoldertitle}\\[4pt]
    {\small feuillets \grothFirst--\grothLast\ (lot \grothBatch)}\\[2pt]
    {\footnotesize\color{gray!60!black}Transcription établie d'après le fac-similé numérisé
     par l'Université de Montpellier.\\
     Les lectures incertaines sont soulignées, les illisibles notées [\,\ldots\,],
     les ajouts entre crochets.}
  \end{center}
  \vspace{1em}}

\endinput


\folder{15}
\batch{4}
\pages{61}{80}
\dating{1965-[vers 1977]}
\watermark{Édition de démonstration}
\foldertitle{Théorie arithmétique. Théorie de Galois des motifs (1965) : notes manuscites (s.d.), tapuscrit (s.d.).}

\begin{document}

\page{62}
\note{la phrase commence au bas du feuillet 60 (lot 3) : « … si $C$ admet
un gén.\ strict $X$, (tout objet ». La rédaction des feuillets 62 à 77 n'est
pas paginée par lui}

de $C$ est de la forme $P \otimes_{o} X$ ($o = \mathrm{End}(X)$), avec $P$
projectif de t.f.\ sur $o$) sous la condition que ces produits tensoriels
soient raisonnables \ill{} ; il suffit \struck{\ill{}} que les
$P \otimes_{o} X$ soient \uncertain{définis}, et indiquer que $C$ est
équivalente à une sous-catégorie pleine de $\mathbb{P}_{d}(o)$
(\uncertain{$o$-modules} à droite projectifs de t.f.) \struck{\ill{}}
contenant l'objet $1$ (et inversement ?). Alors la condition de
\uncertain{plénitude} est automatiquement satisfaite.
\note{« à droite » est ajouté dans l'interligne}

D'autre part, \uncertain{si} \struck{\ill{}} catégorie \struck{\ill{}} $C$
jolie, considérons \struck{les \ill{} objets $X$ de $C$ \ill{}} les
sous-catégories « strictement monogènes » de $C$, celles-ci forment un
\ill{} \uncertain{croissant} de sous-catégories pleines,
\struck{et $C$ \ill{}} et $C$ en est la limite inductive.
\note{« de sous-catégories pleines, » est ajouté dans l'interligne ; le mot
qui précède « croissant » n'est pas lu}

Catégorie \struck{\ill{}} $A$-linéaire \struck{\ill{}} jolie
$\Longleftrightarrow$ limite inductive de sous-catégories pleines
équivalentes chacune \struck{\ill{}} à des catégories \uncertain{du type}
$\mathbb{P}_{d}(o)$, contenant l'\uncertain{élément} $o$,
\note{la page s'arrête sur cette virgule}

\page{65}
$C$ catégorie de Deuring (Picard)

$C \simeq \mathbb{P}(o/A)$

$\mathrm{Brdt}\, C$ opère sur $C$ donc

$\mathrm{Brdt}(C) = \pi_{0}(\mathrm{Brdt}\, C)$ opère sur $\mathrm{Is}(C)$

Le quotient est l'ens.\ des classes d'isomorphie d'algèbres $o$ sur $A$
apparentées à $C$ (i.e.\ à $o$)

Le groupe de stabilité d'un élément $X$, si $\mathrm{Brdt}(A) = 0$, est
isomorphe au groupe $\mathrm{Ext}(o)$ des autom.\ extérieurs de $o$
\note{« $\mathrm{Brdt}(A)$ » : le premier mot est surchargé ; on attendrait
$\mathrm{Pic}(A)$, que la page ne porte pas}

\[
  0 \to o^{Z*} \to o^{*} \to \mathrm{Aut}_{A}\, o \to \mathrm{Brdt}(C)
  \to \mathrm{Is}\, C \to \mathrm{Is}\, \mathcal{O} \to 0
\]
\note{sous les termes de la suite, reliés par des signes d'égalité
verticaux : sous $o^{Z*}$, « $A^{*}$ dans le \uncertain{centre} de $o$ »,
puis $\pi_{1}\, \mathrm{Brdt}(C)$ ; sous $o^{*}$, $\mathrm{Aut}_{C}(X)$,
puis $\pi_{1}(C,X)$ (l'exposant $Z$ semble biffé sur ce second terme) ;
sous $\mathrm{Brdt}(C)$, $\pi_{0}\, \mathrm{Brdt}(C)$ ; sous
$\mathrm{Is}\, C$, $\pi_{0}(C,X)$ ; sous $\mathrm{Is}\, \mathcal{O}$,
« (catégorie des algèbres \uncertain{apparentées} à $o$) » et
« $\pi_{0}(\mathcal{O})$ = \ill{} »}

\begin{itemize}
  \item[] Catégorie de Brandt : $\mathcal{B}$
  \item[] Catégorie de Deuring : $C$
  \item[] Catégorie des anneaux apparentés à $C$ : $\mathcal{O}$
\end{itemize}

$X \in \mathrm{Ob}\, C$, $o = \mathrm{End}(X) \in \mathrm{Ob}\,
\mathcal{O}$, $1 \in \mathrm{Ob}\, \mathcal{B}$
\begin{align*}
  \cdot \to \pi_{1}(\mathcal{B},1) &\to \pi_{1}(C,X) \to
  \pi_{1}(\mathcal{O},o) \to \pi_{0}(\mathcal{B},1) \\
  &\to \pi_{0}(C,X) \to \pi_{0}(\mathcal{O},o) \to 1
\end{align*}
\note{suite encadrée par lui, sur une seule ligne}

\marginal{\struck{\ill{}}}
\note{à gauche de l'encadré, trois lignes écrites de biais, hachurées}

On a envie de dire que $C$ est fibrée sur $\mathcal{O}$, de
\struck{groupe} fibre $\mathcal{B}$
\note{un double trait vertical dans la marge gauche signale ces deux
lignes}

\struck{$C \mapsto A$, $A$ = centre} $A$ est le centre de $o$, …

$\mathrm{Pic}(A) \to \mathcal{B}$ \uncertain{(pl.\ fidèle)}

$\pi_{1}(\mathcal{B},1) = \pi_{1}(\mathrm{Pic}(A)) \simeq A^{*}$

\[
  \to \mathrm{Pic}(A) \to \pi_{0}(\mathcal{B},1) \to \mathrm{Ram}(C/A) \to 0
\]
sous $\mathrm{Ram}(C/A)$ : « groupe de ramification de $C$ sur $A$ »
\[
  \mathrm{Ram}(C/A) \simeq \coprod_{p} \mathbb{Z}/e_{p}\mathbb{Z}
  \quad \text{dans le cas Deuring-Brandt}
\]
\marginal{$\to$ ordres \uncertain{dans} \ill{} \struck{\ill{} des complétés
\ill{}}}
\marginal{\struck{\ill{}} \uncertain{Somme des}
$\mathrm{Br}(\hat{K}_{p})/\mathrm{Br}(\hat{A}_{p})$ ???}
\note{ces deux dernières notes sont au coin inférieur droit, la seconde
entourée ; au-dessus, un premier essai biffé
$\mathrm{Br}(\hat{A}_{p})/p$}

\page{67}
$C$ catégorie de Deuring. On a donc la gr-catégorie $\mathrm{Eqv}(C)$, qui
opère sur $C$, donc
\[
  \mathrm{Brd}(C) = \pi_{0}(\mathrm{Eqv}) \quad \text{opère sur}
  \quad \text{Cl. d'isom. de } C
\]
\note{formule encadrée ; « Classes d'isom.\ de $C$ » est d'abord écrit,
puis biffé}

Ceci dit, le quotient est l'ens.\ des classes d'isomorphie des
$A$-algèbres apparentées à $C$ (donc, dans le cas des \struck{$R$} ordres
maximaux …, le quotient est formé des classes de conjugaison de tous
ordres maximaux), tandis que $\pi_{0}(\mathrm{Eqv})$ correspond
\struck{aux} est le groupe des classes d'isomorphie de $o$-bimodules
= bimodules inversibles de $o$ mod ceux provenant de \struck{Pic} $A^{*}$ ?,
\struck{Notons} que $\mathrm{Is}(C)$ est l'ens.\ des classes d'isomorphie
de $o$-modules : \uncertain{droits} « inv. » \struck{\ill{}} ou objets des
$o$-modules inversibles à droite]. Dans le cas les $o$ centraux
\struck{sur $A$} [on dit $C$ central sur $A$], on a un foncteur pleinement
fidèle
\[
  \mathrm{Pic}(A) \longrightarrow \mathrm{Eqv}(C), \qquad
  L \longmapsto (X \mapsto L \otimes_{A} X)
\]
d'où une suite exacte
\[
  \cdot \to \mathrm{Pic}(A) \to \mathrm{Brdt}(C) \to
  \frac{\mathrm{Brdt}(C)}{\mathrm{Pic}(A)} \to 0
\]
Dans le cas de Brandt (et ordres maximaux, $o$)
\[
  \frac{\mathrm{Brd}\, C}{\mathrm{Pic}(A)} \simeq
  \frac{\mathrm{Div}(C)}{\mathrm{Div}(A)} =
  \coprod_{\substack{p \in \mathrm{Max}\, A \\ p\ \text{ramifié en}\ C}}
  \mathbb{Z}/e_{p}\mathbb{Z}
\]
\note{sous $\mathbb{Z}/e_{p}\mathbb{Z}$ : « (ramification en $p$ »}

\page{69}
Supposons $A$ \struck{local} Dedekind \struck{\ill{}} local (cas Brandt)
\uncertain{du} \ill{} fini

$\mathrm{Is}(C) = 1$, $\mathrm{Brd}(C) = \mathbb{Z}/e\mathbb{Z}$
(ramification

$\mathrm{Pic}(A) = 0$

et on trouve
\[
  \mathrm{Autext}(o) = \mathrm{Brd}\, C = \mathbb{Z}/e\mathbb{Z}
\]
Donc il y a obstruction locale à ce que \uncertain{ses} autom.\ de $o$
soient nec.\ intérieurs, il y a $e$ classes d'autom.\ ext.\ de $o$ …

Supposons (que un autom.\ de $o$ soit localement \uncertain{dans un global}
intérieur (par ex.\ \ill{} \ill{}) [i.e.\ \ill{} dans
$\mathrm{Ram}(C/A)$ nulle]). \struck{Alors les} Il définit alors une
section de $o^{*}/\mathbb{G}_{m,A}$, et \struck{\ill{}} l'obstruction
\ill{} est une classe de $\mathrm{Pic}(A)$, \ill{} \struck{\ill{}}
$H^{1}(A, o^{*})$ i.e.\ donnant une classe triviale dans $\mathrm{Is}(C)$
… \struck{\ill{}} ce qui se voyait aussi par les suites exactes
précédentes.
\note{lecture très incertaine de tout ce paragraphe, surchargé d'ajouts
interlinéaires}

Cas des corps gauches de Deuring \uncertain{ramifiés} en $p$ et $\infty$ :
la suite exacte suivante
\[
  \cdot \to \mathbb{Z}/2\mathbb{Z} \to o^{*} \to \mathrm{Aut}(o) \to
  \mathbb{Z}/2\mathbb{Z} \to \pi_{0}(C,X) \to \pi_{0}(\mathcal{O},o) \to 1
\]
\begin{itemize}
  \item[] sous le premier $\mathbb{Z}/2\mathbb{Z}$ : engendré par $-1$ ;
  \item[] sous $o^{*}$ : $= \mathbb{Z}/2\mathbb{Z}$ sauf si $j = 0$ ou
  $j = 12^{3}$ (où on trouve $\mathbb{Z}/4\mathbb{Z}$ resp.\
  $\mathbb{Z}/6\mathbb{Z}$ si $p \neq 2,3$) \ill{} ;
  \item[] sous le second $\mathbb{Z}/2\mathbb{Z}$ : engendré par
  Frobenius ; et, relié par un trait à la flèche
  $\mathrm{Aut}(o) \to \mathbb{Z}/2\mathbb{Z}$ : épim.\ ssi la courbe
  elliptique $X$ d'anneau $o$ se descend à $\mathbb{F}_{p}$, i.e.\ ssi
  $j \in \mathbb{F}_{p}$ ;
  \item[] sous $\pi_{0}(C,X)$ : classes d'iso.\ des classes de courbes
  elliptiques de Hasse nul \uncertain{mod} l'opération de « Frobenius »
  donnée par $E \mapsto E^{(p)}$ ;
  \item[] sous $\pi_{0}(\mathcal{O},o)$ : classes de conjugaison possibles
  d'ordres maximaux.
\end{itemize}
\note{« resp. » apparie $j = 0$ à $\mathbb{Z}/4\mathbb{Z}$ et $j = 12^{3}$
à $\mathbb{Z}/6\mathbb{Z}$ ; c'est l'appariement inverse ($j = 12^{3}$,
$\mathbb{Z}/4\mathbb{Z}$ ; $j = 0$, $\mathbb{Z}/6\mathbb{Z}$) qu'on
attendrait. Transcrit tel qu'écrit}
\marginal{on voudrait \uncertain{savoir} pour quels $p$ ces $j$ sont
d'inv.\ de Hasse nul …}
\note{note en bas à gauche, reliée par un trait à « $j = 0$ ou
$j = 12^{3}$ » ; voir les calculs du feuillet 76}

Question : $o^{*}$ est-il tjs égal
\note{la phrase s'arrête là ; le feuillet 71 commence autre chose}

\page{71}
Soit \struck{\ill{}} $\mathcal{O}$ la catégorie des $A$-algèbres
apparentées \struck{\ill{}} à une $A$-algèbre déterminée,
\struck{et \ill{} soit} soit $K$ l'anneau total des fractions de $A$,
\struck{considérons $o \in \mathcal{O}$, supposons qu'il existe une partie
$S$ de $A^{\mathrm{reg}}$ \ill{}} \uncertain{dans $A \subset K$}, et
supposons $K$ ½ local.
\note{passage très retravaillé ; les ajouts interlinéaires sont insérés à
la place qu'indiquent ses traits}

Donc \struck{\ill{}} pour $o \in \mathcal{O}$, $o' \in \mathcal{O}$, les
bimodules \ill{} sur $o, o'$ loc.\ libres de rang $1$ \uncertain{des deux
côtés} loc.\ sur \uncertain{$S$-$\mathrm{Spec}\, A$} et libres
\uncertain{isomorphes} sur $o_{K} \simeq E$ et sur $o'_{K}$, dès que
\uncertain{\ill{}} [si tout autom.\ de $E$ est intérieur]
\note{la condition entre crochets est encadrée et soulignée trois fois}

a) $o'_{K} \simeq E$ \struck{\ill{}} \uncertain{donc $o'$ se plonge dans
$E$}, et b) \ill{} $L$ \ill{} avec la structure \ill{} $(o,o')$-bimodule,
les \ill{} plongés dans $E$ comme \ill{}
\marginal{$E$ à isom.\ près \uncertain{un} défini par $o$, seulement
\ill{} de $\mathcal{O}$}

La 2-catégorie \uncertain{$\mathcal{B}$} (pleine \uncertain{de}
$(\mathrm{Cat})$) dont les objets sont les \uncertain{cat.\ de}
\struck{\ill{}} \ill{} : celles de Deuring de type $o$, \struck{et} qui
est équivalente : celle $\mathcal{B}_{1}$ dont les objets sont les objets
de $\mathcal{O}$, $\mathrm{Hom}(o,o')$ étant les catégories des bimodules
sur $o, o'$ qui sont loc.\ libres de rang $1$ \struck{\ill{}} loc.\ sur
$S$, et la composition \ill{} étant
\[
  ({}_{o}L_{o'}, {}_{o'}M_{o''}) \longmapsto L \otimes_{o'} M ,
\]
\note{au-dessus, une première écriture de la même formule, biffée ;
$\mathcal{B}_{1}$ est écrit au bout de la ligne « … : celles »}

$\mathcal{B}_{2}$ est aussi équivalente au 2-groupoïde de Brandt dont les
objets sont les sous-anneaux $o$ de $E$ \struck{apparentés}, $\in
\mathcal{O}$, \struck{\ill{}}
\begin{itemize}
  \item[] $\mathrm{Hom}(o,o') =$ \struck{ensemble des catégorie des}
  $(o,o')$-bimodules $L$ dans $E$ tels que \ill{} \uncertain{des deux
  côtés}, loc.\ libres de rang $1$ loc.\ sur $S$
  \item[] $\mathrm{Hom}(L,M) =$ ens.\ des $\lambda \in E^{\natural}$ tels
  que $\lambda L = M$, si $L, M$ sont des $(o,o')$ bimodules
  « \uncertain{spéciaux} »
\end{itemize}
\note{sous $E^{\natural}$ : « centre de $E$ » ; l'exposant est un petit
signe noté ici $\natural$}

On \uncertain{notera} que les 2-foncteurs $\mathcal{B}_{2} \to
\mathcal{B}_{1} \to \mathcal{B} = \mathcal{B}.$ \uncertain{ne sont pas}
des équivalences « naturelles » !

\page{73}
Supposons maintenant la condition suivante satisfaite (un classique :
\uncertain{tout} autom.\ de $E$ est intérieur)

(C) Pour tout $o \subset E$, $o \in \mathcal{O}$, la loi de composition
dans $\mathrm{Hom}_{\mathcal{B}_{1}}(o,o)$ = ensemble des
$(o,o)$-\struck{biinversibles} spéciaux
\[
  (L, M \mapsto L.M) \quad \text{est commutative :} \quad LM = ML .
\]

\marginal{Je \uncertain{crois} (et il \uncertain{faut si} $A$ est
\uncertain{Dedekind} de dim $1$) que ce sont vrai \uncertain{pour} les
localisés \struck{\ill{}} \uncertain{en idéaux maximaux} (ou leurs
complétés, \uncertain{cf.\ Auslander}). Dans le cas local
\struck{\ill{}} est équivalent à la relation $LM = ML$ pour $L, M$ des
\uncertain{classes} \ill{} idéaux bilatères biinversibles de $o$}

Soit alors $\mathcal{B}_{0}$ une catégorie de Deuring, $C \in
\mathrm{Ob}\, \mathcal{B}_{0}$, et $X \in \mathrm{Ob}\, C$, on pose
\uncertain{et identifie} $o = \mathrm{End}_{C}(X) \in \mathcal{O}$ ; on
peut \uncertain{choisir} \struck{supposer} $E \simeq o_{K}$. Nous
\uncertain{allons} définir sur la catégorie
\[
  \mathrm{Eqv}(C) = \mathrm{Hom}_{\mathcal{B}_{0}}(C,C)
\]
qui forme une cat.\ de Picard stricte \struck{\ill{}}, une
\struck{\ill{}} contrainte de commutativité, \struck{\ill{}}
indépendante du choix de $X$. En effet, on a une équivalence
\[
  \mathrm{Hom}_{\mathcal{B}_{0}}(C,C) \overset{\alpha_{X}}{\simeq}
  \mathrm{Hom}_{\mathcal{B}_{1}}(o,o) = \text{catégorie des }
  (o,o)\text{-bimodules spéciaux}
\]
\note{« spéciaux sous \ill{} » est d'abord écrit, puis biffé et remplacé}
cette dernière catégorie étant cat.\ de Picard stricte comm.\ (en fait,
munie d'une contrainte d'associativité et de commutativité qui sont des
identités). On trouve donc une contrainte de comm.\ stricte sur
$\mathrm{Eqv}(C)$ par transport de structure. Voici l'explicitation
\uncertain{ainsi} : si $L, M$ sont des $(o,o)$-bimodules spéciaux, on
définit
\[
  c : L \otimes_{o} M \simeq M \otimes_{o} L
\]
\struck{\ill{}} en utilisant des isom.\ de $(o,o)$ bimodules
\[
  L \simeq L' \subset E , \qquad M \simeq M' \subset E
\]
avec $L', M'$ des $(o,o)$-bimodules, et notant que
\[
  L \otimes M \simeq L'M' = M'L' \simeq M \otimes L .
\]
\note{les deux dernières égalités sont écrites l'une sous l'autre,
$L'M'$ et $M'L'$ reliés par un signe d'égalité vertical}

\page{75}
Il est immédiat que \struck{cette \ill{}} isom.\ ne dépend pas du choix
des plongements $L \hookrightarrow E$, $M \hookrightarrow E$.

\emph{Proposition.} La contrainte de commutativité précédente ne dépend
pas du choix de $X \in \mathrm{Ob}\, C$.

[\struck{Considérons en effet un autre $X'$ de} Il est clair qu'elle ne
change pas par remplacement de $X$ par un objet isomorphe ; il faudrait
une petite démonstration : dans une gr-catégorie $\mathcal{B}$, les
\struck{foncteurs} autoéquivalences \struck{qui} transforment une
contrainte de comm.\ (stricte ou non) en elle-même …]
\marginal{$\pi_{0}$}

\struck{\ill{}} Soit $X'$ un autre choix, défini par un $P \in
\mathbb{P}(o/A)$. On peut supposer $P \subset E$, soit $o' =
\mathrm{End}(P)$, de sorte que $o'$ est maintenant plongé dans $E$ (c'est
l'ensemble des $\lambda \in E$ tels que $\lambda P \subset P$). Ainsi $P$
est un \struck{bimod} $(o',o)$-bimodule. On a un diagramme commutatif

\begin{tikzcd}[column sep=small, row sep=large, nodes={font=\small}]
 & \mathrm{Eqv}(C,C) \arrow[dl, "\alpha_{X}"'] \arrow[dr] & \\
 (o,o)\text{-bimodules spéciaux} \arrow[rr, "L \mapsto PLP^{-1}"] & & (o',o')\text{-bimodules spéciaux}
\end{tikzcd}

\note{les deux sommets inférieurs portent « Catégorie des … » ; la flèche
de droite n'a pas d'étiquette}

Il suffit donc de prouver que $L \mapsto P.L$ \struck{\ill{}} respecte les
\struck{\ill{}} contraintes de commutativité
\[
  \begin{array}{ccc}
    LM & \text{transformé par } P \text{ en} &
    P(LM)P^{-1} = (PLP^{-1})(PMP^{-1}) \\
    c_{X} = 1 \ \| & & \| \ P(c_{X}) = 1 \qquad \| \ c_{X'} = 1 \\
    ML & & P(M.L)P^{-1} = (PMP^{-1})(PLP^{-1})
  \end{array}
\]

OK.

\page{76}
\note{notes éparses de sa main, écrites en travers d'un feuillet du
cours dactylographié étranger à ce dossier ; transcrites sans ordre,
dans leur disposition de haut en bas}

$\mathbb{Z}_{p}($ ; $\mathbb{Q}(i)^{*}$ ; $T$

$\mathcal{B}_{c}(\mathbb{Z})$

$-1$ carré mod $p$ ; $p \equiv 1$ $(4)$

$\left[\dfrac{-1}{p}\right] = -1$

\uncertain{TORSC}
\note{près de « $-1$ », une formule entourée, illisible ; ces calculs
semblent répondre à la question du feuillet 69 (pour quels $p$ les
valeurs $j = 0$, $j = 12^{3}$ sont d'invariant de Hasse nul)}

\page{77}
\emph{Conclusion.} Les catégories $\mathrm{Eqv}(C)$, $C \in \mathrm{Ob}\,
\mathcal{B}_{0}$, sont munies canoniquement d'une structure de
\struck{des gr-catégories munies de str.} catégorie de Picard strictes
\struck{commutatives}. \struck{Elles sont} Ces cat.\ de Picard sont
canoniquement équivalentes entre elles, et peuvent donc s'identifier à une
\struck{\ill{}} catégorie de Picard $\mathcal{B}$ (catégorie de Picard de
Brandt). Donc $\mathcal{B}$ opère sur chaque $C$, et pour $C, C' \in
\mathrm{Ob}\, \mathcal{B}_{0}$, la catégorie \struck{Hom}
$\mathrm{Eqv}(C,C')$ est canoniquement munie d'une structure de
torseur-catégorie sous $\mathcal{B}$. De façon précise,
\[
  \begin{array}{ccc}
    & C' \longmapsto \mathrm{Eqv}(C,C') & \\
    \mathcal{B}_{0} & \longrightarrow &
    \text{la catégorie des torseurs-catégories sous } \mathcal{B}
  \end{array}
\]
est une équivalence de 2-catégories.

Si on a un chgt de base $A \to A'$, et si les $C_{A'}$ satisfait les mêmes conditions que les $C$, alors on a un homom.\
de \struck{gr-}catégories de Picard
\[
  \mathcal{B} \longrightarrow \mathcal{B}'
\]
et l'homom.\ \struck{\ill{}}
\[
  C \longrightarrow C_{A'}
\]
est compatible avec les opérations de $\mathcal{B}$, $\mathcal{B}'$ …
\note{sous $A \to A'$, un petit carré : $A \subset K$, $A' \subset K'$,
$K \to K'$}

\section{Localisation pour les variétés abéliennes}

\note{titre dactylographié, en tête du feuillet 80. Le feuillet est le
premier d'un tapuscrit ; les ajouts manuscrits sont donnés comme notes
marginales de sa main, les corrections faites à la machine sont
transcrites directement. Comme souvent sur ces frappes, les flèches manquent ;
elles sont rétablies et signalées}

\page{80}
1. Soit $A$ un anneau commutatif. Soit $C$ une catégorie $A$-linéaire (pas
nécessairement additive ; les Hom sont munis de structures $A$-linéaires,
les accouplements sont $A$-bilinéaires). Pour toute \struck{\ill{}}
$A$-algèbre $A'$, soit $C_{A'}$ la catégorie ayant mêmes objets que $C$, et
où les Hom sont
\[
  \mathrm{Hom}_{C_{A'}}(X,Y) = \mathrm{Hom}_{C}(X,Y) \otimes_{A} A' .
\]
\marginal{$\underline{C}$}
\note{lettre doublement soulignée, dans la marge gauche, en regard de la
phrase suivante}
On trouve ainsi une \struck{champ en} catégorie fibrée
\marginal{$\underline{C}$} en catégories \struck{$A$}-linéaires sur le site
annelé des schémas affines $\mathrm{Spec}(A')$ sur $S = \mathrm{Spec}(A)$.
Elle n'est pratiquement jamais un champ sur ce gros site, même pour une
topologie assez grossière comme la topologie de Zariski. Cependant, si on
se restreint au petit site \add{affine} de Zariski de $S$ des ouverts
affines de $S$, on a souvent un champ (pour la topologie de Zariski), soit
$\underline{C}$. \struck{Ceci a des conséquences}
\marginal{On dit alors que $C$ est localisable sur $A$.}
Supposons cette condition satisfaite. On a des conséquences agréables,
comme :
\note{« catégorie fibrée » et « affine » sont frappés dans l'interligne ;
« soit $\underline{C}$ » aussi}

a) Soit $R$ une $A$-algèbre (pas néc.\ commutative), et $X$ un objet de
$C$ muni d'une opération $A$-linéaire de $R$. Alors pour tout $R$-module
à droite $P$, qui est isomorphe localement sur $S$ à $R_{d}$, le produit
tensoriel $P \otimes_{R} X$ dans $C$ existe, et si $R \to
\mathrm{End}_{C}(X)$ est bijectif, le foncteur
\[
  \underline{P}(R) \longrightarrow C , \qquad P \longmapsto P \otimes_{R} X
\]
de la catégorie des $R$-modules à droite localement isomorphes à $R$ sur
$S$ dans la catégorie $C$ est pleinement fidèle,
\marginal{et son image essentielle est formée des $X \in \mathrm{Ob}\, C$
qui sont « localement isomorphes à $X$ » sur $A$.}
\note{« si $R \to \mathrm{End}_{C}(X)$ est bijectif » est frappé dans
l'interligne, sans la flèche ; dans le foncteur, la flèche manque aussi et
$\underline{P}(R)$ est frappé sous la place où elle devait aller}

b) Supposons que $C$ soit additive. Alors $\underline{C}$ est un champ en
catégories additives $A$-linéaires. Pour tout $X$, $R$ comme dessus,
\struck{et} et tout module à droite $P$ sur $R$ qui est localement sur $S$
libre de type fini, le produit tensoriel $P \otimes_{R} X$ existe, et si
$R \to \mathrm{End}_{C}(X)$, le foncteur
\[
  P \longmapsto P \otimes_{R} X
\]
\note{la frappe s'arrête ici, au bas du feuillet ; la phrase continue au
feuillet 81 (lot 5), « de la catégorie de ces modules vers $C$ est
pleinement fidèle, … »}

\end{document}
